% Generated by Listas/make_statements.py from pratica-lecture-02.md, item 7. Do not edit.
\textbf{Source.} Tadelis (2013), \emph{Game Theory: An Introduction}, Exercise 6.1 ``Best Responses in the Battle of the Sexes'', p.~123 {[}PDF p.~142{]}. Answer in the source: none (the preface, p.~xiii, says about half of the exercises are solved on Princeton UP's website). Reproduced verbatim.



\begin{quote}

\textbf{Best Responses in the Battle of the Sexes:} Use the best-response correspondences in the Battle of the Sexes game to find all the Nash equilibria. (Follow the approach used for the example in Section 6.2.3.)

\end{quote}



\textbf{Referred to in the statement.} The Battle of the Sexes (Tadelis §3.3, pp.~53--54 {[}PDF p.~72--73{]}) (reproduced from the book):



\begin{quote}

Alex and Chris are a couple, and they need to choose where to meet this evening, with no means of communicating. Both prefer being together over not being together, but Alex prefers opera (\(O\)) to football (\(F\)), while Chris prefers the opposite.



{\def\LTcaptype{none} % do not increment counter

\begin{longtable}[]{@{}lll@{}}

\toprule\noalign{}

Alex ~Chris & \(O\) & \(F\) \\

\midrule\noalign{}

\endhead

\bottomrule\noalign{}

\endlastfoot

\(O\) & 2, 1 & 0, 0 \\

\(F\) & 0, 0 & 1, 2 \\

\end{longtable}

}

\end{quote}

